\textbf{Phần 1. Trắc nghiệm lựa chọn}
\setcounterpageref{ex}{0}

\begin{ex} %[2D13V67]
Trong giờ thực hành, bạn Nam ghi lại thời tiết trong 10 ngày: có 6 ngày nắng, 3 ngày mưa và 1 ngày râm mát. Ngày có thời tiết xuất hiện nhiều nhất là
\choice
{\True Ngày nắng}
{Ngày mưa}
{Ngày râm mát}
{Cả nắng và mưa bằng nhau}
\loigiai{
Ngày nắng xuất hiện 6 ngày, nhiều nhất.
\textbf{Đáp án của câu hỏi là A}
}
\end{ex}

\begin{ex} %[2D13V67]
Khảo sát loại quả yêu thích của 30 học sinh: 12 bạn thích xoài, 10 bạn thích cam, 8 bạn thích dưa hấu. Số bạn thích dưa hấu ít hơn số bạn thích xoài là
\choice
{$2$ bạn}
{$6$ bạn}
{\True $4$ bạn}
{$5$ bạn}
\loigiai{
$12 - 8 = 4$ (bạn).
\textbf{Đáp án của câu hỏi là C}
}
\end{ex}

\begin{ex} %[2D13V67]
Để đếm số chim bồ câu đậu trên mái nhà, cách kiểm đếm nào sau đây là nhanh và không bị nhầm lẫn?
\choice
{Đếm nhảy cóc tùy ý}
{\True Đếm lần lượt từng con từ trái sang phải rồi gạch ghi lại}
{Ước lượng đại khái không cần đếm}
{Nhắm mắt đoán số lượng}
\loigiai{
Đếm có trật tự từ trái sang phải và ghi lại là cách chuẩn khoa học.
\textbf{Đáp án của câu hỏi là B}
}
\end{ex}

\begin{ex} %[2D13V67]
Khi tung đồng xu 10 lần, bạn Mai ghi nhận được 6 lần xuất hiện mặt Sấp. Hỏi có bao nhiêu lần xuất hiện mặt Ngửa?
\choice
{$5$ lần}
{$3$ lần}
{$6$ lần}
{\True $4$ lần}
\loigiai{
$10 - 6 = 4$ (lần).
\textbf{Đáp án của câu hỏi là D}
}
\end{ex}

\begin{ex} %[2D13V67]
Bảng thống kê số cây mỗi bạn trồng được: An trồng 3 cây, Bình trồng 5 cây, Chi trồng 4 cây, Dũng trồng 2 cây. Bạn trồng được nhiều cây nhất là
\choice
{\True Bình}
{An}
{Chi}
{Dũng}
\loigiai{
Bình trồng 5 cây, nhiều nhất.
\textbf{Đáp án của câu hỏi là A}
}
\end{ex}

\begin{ex} %[2D13V67]
Từ bảng số liệu trên, tổng số cây cả bốn bạn An, Bình, Chi, Dũng trồng được là
\choice
{$12$ cây}
{$15$ cây}
{\True $14$ cây}
{$16$ cây}
\loigiai{
$3 + 5 + 4 + 2 = 14$ (cây).
\textbf{Đáp án của câu hỏi là C}
}
\end{ex}

\begin{ex} %[2D13V67]
Thực hành cân đo sức khỏe: Bạn Hùng nặng $28\text{ kg}$, bạn Tuấn nặng $25\text{ kg}$, bạn Dũng nặng $31\text{ kg}$. Bạn nặng cân nhất là
\choice
{Bạn Hùng}
{\True Bạn Dũng}
{Bạn Tuấn}
{Hai bạn nặng bằng nhau}
\loigiai{
Bạn Dũng nặng 31 kg, lớn nhất.
\textbf{Đáp án của câu hỏi là B}
}
\end{ex}

\begin{ex} %[2D13V67]
Trong hộp có 10 tấm thẻ ghi các số từ 1 đến 10. Rút ngẫu nhiên một tấm thẻ. Sự kiện rút được tấm thẻ ghi số có hai chữ số là
\choice
{Chắc chắn xảy ra}
{Không thể xảy ra}
{Luôn luôn xảy ra}
{\True Có thể xảy ra}
\loigiai{
Trong các số từ 1 đến 10 có số 10 là số có hai chữ số nên có thể xảy ra.
\textbf{Đáp án của câu hỏi là D}
}
\end{ex}

\begin{ex} %[2D13V67]
Bạn Nam gieo con xúc xắc 20 lần và ghi lại số lần xuất hiện mặt 6 chấm được 3 lần. Số lần xuất hiện các mặt khác là
\choice
{\True $17$ lần}
{$16$ lần}
{$18$ lần}
{$15$ lần}
\loigiai{
$20 - 3 = 17$ (lần).
\textbf{Đáp án của câu hỏi là A}
}
\end{ex}

\begin{ex} %[2D13V67]
Theo dõi số học sinh đi muộn trong tuần của lớp: Thứ Hai có 2 bạn, thứ Ba có 0 bạn, thứ Tư có 1 bạn, thứ Năm có 0 bạn, thứ Sáu có 3 bạn. Có bao nhiêu ngày trong tuần không có bạn nào đi muộn?
\choice
{$1$ ngày}
{$3$ ngày}
{\True $2$ ngày}
{$0$ ngày}
\loigiai{
Đó là thứ Ba và thứ Năm (2 ngày).
\textbf{Đáp án của câu hỏi là C}
}
\end{ex}

\begin{ex} %[2D13V67]
Tổng số lượt học sinh đi muộn trong cả tuần trên là
\choice
{$5$ lượt}
{\True $6$ lượt}
{$7$ lượt}
{$4$ lượt}
\loigiai{
$2 + 0 + 1 + 0 + 3 = 6$ (lượt).
\textbf{Đáp án của câu hỏi là B}
}
\end{ex}

\begin{ex} %[2D13V67]
Khi thu thập thông tin về chiều cao của các bạn trong tổ, dụng cụ đo thích hợp nhất là
\choice
{Cân đồng hồ}
{Đồng hồ bấm giây}
{Nhiệt kế y tế}
{\True Thước đo chiều cao (thước dây hoặc thước cuộn)}
\loigiai{
Đo chiều cao phải dùng thước dây hoặc thước cuộn.
\textbf{Đáp án của câu hỏi là D}
}
\end{ex}

\textbf{Phần 2. Trắc nghiệm đúng sai}
\setcounterpageref{ex}{0}

\begin{ex} %[2D13V67]
Xét tính đúng sai trong thực hành thu thập và xử lý số liệu:
\choiceTF
{\True Ghi chép số liệu cần trung thực, chính xác và rõ ràng.}
{Chỉ cần phán đoán theo cảm tính mà không cần đếm hay đo thực tế.}
{\True Bảng thống kê giúp ta dễ dàng so sánh các đại lượng và rút ra kết luận.}
{\True Sau khi thu thập, việc kiểm đếm lại giúp tránh nhầm lẫn.}
\loigiai{
\begin{itemize}
    \item Ý a) Đúng: Nguyên tắc khoa học. Đúng.
    \item Ý b) Sai: Khoa học thực nghiệm đòi hỏi đo đếm thực tế. Sai.
    \item Ý c) Đúng: Đúng vai trò của thống kê.
    \item Ý d) Đúng: Đúng.
\end{itemize}
\begin{itemize}
    \item a) \textbf{Đúng}.
    \item b) \textbf{Sai}.
    \item c) \textbf{Đúng}.
    \item d) \textbf{Đúng}.
\end{itemize}
}
\end{ex}

\begin{ex} %[2D13V67]
Khảo sát màu sắc yêu thích của 20 bạn học sinh: Màu xanh có 8 bạn, màu đỏ có 7 bạn, màu vàng có 5 bạn:
\choiceTF
{\True Màu xanh được nhiều bạn yêu thích nhất.}
{\True Số bạn thích màu vàng ít hơn màu đỏ là 2 bạn.}
{Tổng số bạn tham gia khảo sát là 25 bạn.}
{\True Số bạn thích màu xanh hoặc màu vàng là 13 bạn.}
\loigiai{
\begin{itemize}
    \item Ý a) Đúng: 8 > 7 > 5. Đúng.
    \item Ý b) Đúng: 7 - 5 = 2 bạn. Đúng.
    \item Ý c) Sai: 8 + 7 + 5 = 20 bạn, không phải 25. Sai.
    \item Ý d) Đúng: 8 + 5 = 13 bạn. Đúng.
\end{itemize}
\begin{itemize}
    \item a) \textbf{Đúng}.
    \item b) \textbf{Đúng}.
    \item c) \textbf{Sai}.
    \item d) \textbf{Đúng}.
\end{itemize}
}
\end{ex}

\textbf{Phần 3. Trả lời ngắn}
\setcounterpageref{ex}{0}

\begin{ex} %[2D13V67]
Khảo sát sở thích đọc sách của 40 học sinh: 18 bạn thích truyện cổ tích, 14 bạn thích truyện tranh, còn lại thích sách khoa học. Có bao nhiêu bạn thích sách khoa học?
\shortans{8}
\loigiai{
\textbf{Đáp số:} $8$\\
$40 - 18 - 14 = 8$ bạn.
}
\end{ex}

\begin{ex} %[2D13V67]
Tung đồng xu 15 lần, ghi nhận được 8 lần mặt Ngửa. Hỏi có bao nhiêu lần mặt Sấp xuất hiện?
\shortans{7}
\loigiai{
\textbf{Đáp số:} $7$\\
$15 - 8 = 7$ lần.
}
\end{ex}

\begin{ex} %[2D13V67]
Đếm số ô cửa sổ của trường học: tầng 1 có 12 ô, tầng 2 có 14 ô, tầng 3 có 14 ô. Có tất cả bao nhiêu ô cửa sổ?
\shortans{40}
\loigiai{
\textbf{Đáp số:} $40$\\
$12 + 14 + 14 = 40$ ô.
}
\end{ex}

\begin{ex} %[2D13V67]
Trong đợt thi đua, tổ Một đạt 28 bông hoa điểm tốt, tổ Hai đạt nhiều hơn tổ Một 6 bông hoa. Cả hai tổ đạt bao nhiêu bông hoa điểm tốt?
\shortans{62}
\loigiai{
\textbf{Đáp số:} $62$\\
Tổ 2: $28 + 6 = 34$. Cả hai tổ: $28 + 34 = 62$ bông hoa.
}
\end{ex}

\begin{ex} %[2D13V67]
Một nhóm gồm 4 bạn, mỗi bạn đếm được 5 chiếc xe buýt chạy qua. Nhóm đã đếm được tổng cộng bao nhiêu chiếc xe buýt?
\shortans{20}
\loigiai{
\textbf{Đáp số:} $20$\\
$4 \times 5 = 20$ chiếc xe buýt.
}
\end{ex}

\begin{ex} %[2D13V67]
Trong túi có 6 viên bi đỏ và 4 viên bi trắng. Cần lấy ít nhất mấy viên để chắc chắn có 1 viên bi đỏ?
\shortans{5}
\loigiai{
\textbf{Đáp số:} $5$\\
Trường hợp xấu nhất lấy hết 4 bi trắng, lấy thêm 1 viên nữa: $4 + 1 = 5$ viên.
}
\end{ex}
